\chapter{Two Languages, Two Philosophies}
\label{ch:philosophy}

\chapabstract{This chapter answers one question: why does \SV{} feel wrong to
someone who has written \VHDL{} for five years? The answer is not syntax. The
two languages were built for different purposes, by different communities, with
different ideas about where correctness comes from. By the end of the chapter
you will know where \SV{} puts the guard rails that \VHDL{} puts in the type
system, how the \SV{} event scheduler differs from the \VHDL{} simulation
cycle, why blocking and non-blocking assignments exist at all, and how the
verification culture around \SV{} differs from the one around \VHDL{}. The
chapter ends with the other half of the argument: an honest account of what
\VHDL{} does better, and a set of criteria for choosing between the two.}

% ---------------------------------------------------------------------------
% Chapter-local helper: a Verilog compiler directive in running text.
% ---------------------------------------------------------------------------

You already know how to design digital hardware. Nothing in this book will
teach you what a flip-flop is, what a handshake is, or why a long combinational
path costs you frequency. What this book teaches is a second language for
saying things you can already say, plus a set of things that the second
language says which the first one cannot.

The temptation, in the first week, is to treat \SV{} as \VHDL{} with different
punctuation. Look up the keyword, substitute it, move on. That works for about
two days. Then a design simulates correctly and synthesises into something
else, or a testbench passes for the wrong reason, or a signal is one bit wide
when you were certain you declared it as eight. At that point the missing
knowledge is not a keyword. It is the model of the world that the language was
built on. So this chapter is deliberately not a syntax table.
\cref{ch:rosetta} and \cref{ch:rosetta2} are the syntax tables. This chapter
is about where the two languages come from, and which of your instincts have
to be retrained before those tables are safe to use.

The order is: two origin stories, then the type system you are about to lose,
then design units and the preprocessor, then the scheduler, which is the
section that matters most, then declared intent, then the verification culture.
The last two sections turn the argument round. Having spent the chapter on
what \SV{} gives you, they close with what \VHDL{} does better, which is more
than a book with this title is normally willing to admit, and with criteria
for deciding which language a given job wants. Read them: a designer who can
only argue one way is not much use in a language review.

% ===========================================================================
\section{Two Origins, and Why They Still Matter}
\label{sec:philosophy:origins}
% ===========================================================================

Languages carry their origins for decades. \VHDL{} and Verilog were designed in
the same half-decade, on the same continent, for the same industry, and they
still ended up almost as different as two hardware description languages can
be, because the problems their authors were paid to solve were not the same
problem.

\subsection{VHDL: a documentation language that learned to synthesise}

\index{VHDL!history}\VHDL{} began in 1981 as part of the United States
Department of Defense Very High Speed Integrated Circuit (VHSIC) programme. The
problem the DoD had was not simulation speed and it was not designer
productivity. It was procurement. A military system built around a custom chip
might have to be maintained, re-qualified, and possibly re-fabricated twenty or
thirty years after the original vendor had lost interest or gone out of
business. The DoD wanted a vendor-neutral, machine-readable, unambiguous
description of a circuit that a different contractor could read, understand and
re-simulate decades later.

That is a specification for a documentation language that happens to be
executable. Everything about \VHDL{} follows from it. The syntax was derived
from Ada, which was itself the DoD's standard programming language and which
had been designed by committee for long-lived, safety-critical, multi-author
software. \VHDL{} inherited Ada's verbosity, its strong typing, its explicit
declarations, its named association, its packages and libraries, and its
general belief that saying a thing twice in slightly different ways is a
feature, because it lets the compiler check that you meant it.

The language was released publicly in 1985 and standardised as IEEE 1076-1987;
the revisions you care about are IEEE 1076-1993, IEEE 1076-2008, which is what
this book assumes, and IEEE 1076-2019, whose tool support is still
thin.\index{IEEE 1076} Note what is missing from that timeline: synthesis.
\VHDL{} was designed to be simulated, not to be turned into gates. Synthesis
arrived from outside, in the early 1990s, as a commercially interesting subset
that particular tools agreed to support. IEEE 1076.6 tried to standardise that
subset and was eventually withdrawn. To this day the synthesisable part of
\VHDL{} is a convention enforced by tools, not a part of the language
definition.

\subsection{Verilog: a simulator that learned to be a language}

\index{Verilog!history}Verilog has the opposite biography. Around 1984, Gateway
Design Automation, a small company, sold a logic simulator called Verilog-XL.
Verilog was its input format. Phil Moorby designed it, and he designed it for
one purpose: to be fast to write, fast to parse, and fast to simulate on the
hardware of 1984.

The language is therefore C-flavoured, terse, and deeply pragmatic. It has
begin/end instead of Ada's fully spelled-out closers. It has a bit-vector type
and an integer and very little else, because that is what an event simulator
needs. It has four-valued logic with strengths, because a 1984 simulator had to
model pull-ups, tri-states and buses at the transistor level. It has an
explicit textual preprocessor, because C had one and it was cheap to implement.
And most importantly, its semantics are the semantics of its simulator: the
language was not specified first and implemented afterwards, it was implemented
first and specified afterwards, from the behaviour of Verilog-XL.

Cadence acquired Gateway in 1989. Under commercial pressure from \VHDL{}, which
was an open standard, Cadence put Verilog into the public domain in 1990 and
Open Verilog International standardised it; it became IEEE 1364-1995, then
1364-2001 (which added the things RTL designers actually needed, including
ANSI-style port lists, \code{generate}, and signed arithmetic) and 1364-2005.
\index{IEEE 1364}

\subsection{SystemVerilog: a verification language that swallowed its host}

\index{SystemVerilog!history}By the late 1990s Verilog was losing the
verification argument. Testbenches were being written in dedicated verification
languages -- Synopsys Vera, Verisity's \textit{e} -- because Verilog had no
classes, no randomisation and no functional coverage. At the same time Co-Design
Automation was developing Superlog, an extension of Verilog with C-like data
types and higher-level constructs.

Accellera, an industry consortium, assembled \SV{} out of these donated pieces:
Superlog for the design constructs and the type system, Vera for the
object-oriented verification layer, and a separate assertion effort for
SystemVerilog Assertions. \SV{} 3.0 (2002) was a design-side extension; \SV{}
3.1 (2003) added the verification half; Accellera \SV{} 3.1a became IEEE
1800-2005.

The decisive event is IEEE 1800-2009. Until then, \SV{} was formally defined as
a set of extensions \emph{on top of} IEEE 1364 Verilog, and you needed both
documents. In 2009 the two standards were merged: IEEE 1800-2009 contains all
of Verilog, and IEEE 1364 was withdrawn. Verilog is no longer a language with a
standard of its own. It is a subset of \SV{}. The current revisions are IEEE
1800-2017, which this book targets, and IEEE 1800-2023.\index{IEEE 1800}

\begin{notebox}{What ``the LRM'' means}
Engineers say \emph{the \LRM{}} for the Language Reference Manual: IEEE 1800 for
\SV{}, IEEE 1076 for \VHDL{}. IEEE 1800-2017 is about 1300 pages, roughly half
of which describes things you will never synthesise. When this book says ``the
\LRM{} requires'', you can check it: clause 4 is the scheduling semantics,
clause 6 the data types, clause 9 the processes, clause 23 the modules and
hierarchy.
\end{notebox}

\subsection{The practical consequence}

Here is the sentence to remember from this section. \VHDL{} was designed as a
language and later grew a synthesis subset. \SV{} was designed as a simulator
input and later grew a type system.

Both statements have visible consequences in code you will write next week.
Because \VHDL{} was designed as a language first, its rules are uniform,
checkable and occasionally infuriating: the analyser will reject a design that
would have simulated perfectly well, because a rule says so. Because \SV{} grew
outward from a simulator, its rules are permissive by default: almost anything
you type has \emph{some} defined meaning, so the tool accepts it and gives you
that meaning, whether or not it is the one you had in mind. Neither approach is
stupid. \VHDL{} optimises for the case where the reader is not the author and
the year is 2045; \SV{} optimises for the case where the author is at the
terminal now and wants a waveform.

% ===========================================================================
\section{The Compiler Is Not Your Friend Any More}
\label{sec:philosophy:typing}
% ===========================================================================

In \VHDL{} you have an instinct, built over years, that goes like this: if it
analyses and elaborates, the types are right. That instinct is correct in
\VHDL{} and it is worth a great deal. You cannot assign an \code{integer} to a
\code{std\_logic\_vector}, you cannot add \code{1} to one without deciding in
writing whether it is signed or unsigned, and you cannot connect a 12-bit port
to a 16-bit signal. Each of those is a compile-time or elaboration-time
failure, and each corresponds to a real design bug you therefore never
shipped.

\SV{} does not work this way, and the single most dangerous thing a \VHDL{}
designer can do is carry the instinct across unchanged.

\subsection{Widths do not have to match}

\index{width mismatch}In \SV{}, assignment between vectors of different widths
is legal. The right-hand side is truncated from the most significant end if it
is too wide, and extended if it is too narrow: with zeros if the expression is
unsigned, by sign extension if it is signed. No warning is required by the
\LRM{}. \cref{lst:philosophy:width} shows the whole problem in six lines.

\begin{svcode}[caption={Width rules in action. Nothing here is an error.},label={lst:philosophy:width}]
module width_demo;
  logic [7:0] a;
  logic [3:0] b;
  logic [3:0] sum4;
  logic [4:0] sum5;

  initial begin
    a    = 8'hAB;
    b    = a;             // b = 4'hB   : silently truncated
    sum4 = 4'hF + 4'h1;   // sum4 = 4'h0: the carry is lost
    sum5 = 4'hF + 4'h1;   // sum5 = 5'h10: operands widened to 5 bits
    $display("b=%h sum4=%h sum5=%h", b, sum4, sum5);
  end
endmodule
\end{svcode}

Look carefully at the last two assignments, because they are the same
expression producing two different results. \SV{} computes the width of an
arithmetic expression from its \emph{context}, and the assignment target is
part of that context. In line 10 the target is four bits wide, so the addition
is done in four bits and the carry out disappears. In line 11 the target is
five bits wide, so both operands are first extended to five bits and the carry
is kept. The expression did not change. The declaration of the variable you
assigned it to changed the arithmetic.

\VHDL{} has no equivalent behaviour because it has no equivalent hole. The
corresponding code does not compile, and the fix is forced to be explicit.

\begin{rosetta}
\begin{rleft}
signal a    : unsigned(7 downto 0);
signal b    : unsigned(3 downto 0);
signal sum5 : unsigned(4 downto 0);

-- b <= a;  -- length mismatch: rejected
b    <= a(3 downto 0);
sum5 <= resize(a(3 downto 0), 5)
      + to_unsigned(1, 5);
\end{rleft}
\begin{rright}
logic [7:0] a;
logic [3:0] b;
logic [4:0] sum5;

// b = a;   // legal, truncates
assign b    = a[3:0];
assign sum5 = 5'(a[3:0]) + 5'd1;
\end{rright}
\end{rosetta}

The \VHDL{} side says \code{resize} because it must. The \SV{} side does not
have to say anything, which is exactly why you should teach yourself to say
something anyway: the cast \code{5'(...)} documents the intent and makes the
widening visible to the next reader. \cref{ch:rosetta} covers the width and
signedness rules in full, including the difference between \code{>>} and
\code{>>>}, which catches people at least as often.

\begin{gotcha}{Your ports do not have to match either}
Port width mismatch at instantiation is also legal. Connecting a 16-bit signal
to an 8-bit input port silently drops the top eight bits; connecting an 8-bit
signal to a 16-bit port zero-extends it. In \VHDL{} this is an elaboration
error and your synthesis run stops. In \SV{} it is a Tuesday. This is the
single most common cause of ``it worked in the block-level testbench and broke
in the system'' among people new to the language. Named port connections
(\code{.data\_i(my\_data)}) protect you against wiring the wrong port; they do
not protect you against wiring it at the wrong width.
\end{gotcha}

\subsection{Types that look familiar and are not}

\index{logic}\index{std\_logic}\code{logic} is the type you will use for
essentially everything in RTL, and the first thing to say about it is that it
is \emph{not} \code{std\_logic}. It is close: it is a four-valued type with the
values \code{0}, \code{1}, \code{x} and \code{z}. But it has no strengths.
\code{std\_logic} has nine values, including the weak strengths \code{L},
\code{H}, \code{W} and the don't-care \code{-}. Those extra values matter in
two places you have probably used them: pull-up modelling, and \code{-} in a
\code{case} choice or a comparison, which is a trap in \VHDL{} anyway because
\code{-} does not mean don't-care to the \code{=} operator. \SV{} has no
\code{-}. Its don't-care in a case statement is a different mechanism
(\code{casez}, or better, the \code{unique} and \code{priority} qualifiers),
covered in \cref{ch:rosetta2}.

The second thing to say about \code{logic} is subtler and more important.
\code{logic} is not a wire. In original Verilog there were two families,
\emph{nets} (\kw{wire}, \code{tri}, \code{wand}) which model connections and
can have several drivers resolved together, and \emph{variables} (\kw{reg})
which model storage inside a process and may have only one driver.
\code{logic}\index{wire}\index{reg} is simply a better name for \kw{reg}: it is
a single-driver, four-valued variable that you may also drive with a single
continuous assignment. If you genuinely need multiple drivers resolved -- a
tri-state bus -- you need \kw{wire}. The \VHDL{} habit of writing
\code{std\_logic} everywhere and relying on the resolution function translates
to \code{logic} everywhere and \emph{not} relying on resolution, because
\code{logic} will not resolve; a second driver is an error.

Then there are the two-state types, which have no \VHDL{} counterpart at all
and are a genuine trap.

\begin{gotcha}{\code{int} is not \code{integer}, and it hides your reset bugs}
\SV{} has two-state types: \kw{bit} (1 bit), \kw{byte} (8), \kw{shortint} (16),
\kw{int} (32), \kw{longint} (64). Two-state means the variable can only hold
\code{0} or \code{1}. It cannot hold \code{x}. It initialises to \code{0}, not
to \code{x}.

That last sentence is the trap. A \VHDL{} designer relies on a signal starting
at \code{'U'} or \code{'X'} and propagating that unknown through the design
until reset drives it: a red waveform after reset is how you find the register
you forgot to reset. Declare that register as \kw{int} or \kw{bit} and it
starts at zero, the simulation looks clean, and the silicon powers up
arbitrarily. Use \code{logic} for anything that is hardware. Keep the two-state
types for loop counters, testbench bookkeeping and class members, where the
speed and the absence of \code{x} are what you want.

Note also that \SV{}'s \kw{int} is 32-bit and signed by definition, whereas
\VHDL{}'s \code{integer} is a 32-bit signed type whose range is defined by the
standard but whose subtype you are expected to constrain. There is no \SV{}
equivalent of \code{natural} or of \code{range 0 to 15}, and therefore no
runtime range check.
\end{gotcha}

\subsection{Names you never declared}

\index{implicit net declaration}Verilog inherited from its simulator days a
rule that would be unthinkable in \VHDL{}: some undeclared identifiers are not
errors. If a name that has never been declared appears in an instance port
connection, or on the left-hand side of a continuous assignment, the tool
creates a one-bit net called that, and carries on.

\begin{svcode}[caption={A typo that costs you an afternoon},label={lst:philosophy:implicit}]
module counter (input logic clk, output logic [3:0] cnt_o);
  logic [3:0] cnt_d, cnt_q;

  assign cnt_d = cnt_q + 4'd1;

  sink u_sink (.d_i(cnt_qq), .q_o(cnt_o));  // typo: cnt_qq never declared

  always_ff @(posedge clk) cnt_q <= cnt_d;
endmodule
\end{svcode}

\cref{lst:philosophy:implicit} compiles. \code{cnt\_qq} becomes a one-bit
implicit net, driven by nothing, so it is \code{z}. The port it feeds is four
bits wide, so the one-bit net is padded with zeros and the submodule sees
\code{4'b000z}: the top bits are a clean \code{0} and only bit 0 carries the
\code{z}. That is the reason the bug is hard to see. A wholly undriven signal
shows up as \code{zzzz} in the waveform and in every print; \code{000z} looks
almost like data, and any arithmetic on it quietly turns into \code{x}. In
\VHDL{} this is one line of error output naming the undeclared identifier.

The remedy is not optional. Put the directive
\chOneDir{default\_nettype none} at the top of every file: it switches
implicit net creation off, and the typo above becomes the compile error it
should always have been.\index{default\_nettype}

\begin{tipbox}{The two lines that go in every file}
Begin every \file{.sv} file with \chOneDir{default\_nettype none} and end it
with \chOneDir{default\_nettype wire}. The closing line matters because
compiler directives are not scoped to a file: they persist through the whole
compilation, so a file that switches the default off and does not switch it
back will break the next file the tool reads, which may be somebody else's
IP.

\begin{svcode}
`default_nettype none

module fifo #(parameter int unsigned DEPTH = 8) (
  input  wire  logic clk_i,
  input  wire  logic rst_ni,
  // ...
);
  // ...
endmodule

`default_nettype wire
\end{svcode}
Once implicit nets are off, undeclared input ports also become errors, which is
why the ports above are written \code{input wire logic}: with the default net
type set to none you must say what kind of thing an input port is.
\end{tipbox}

\subsection{The linter is the new type checker}

\index{linting}\index{Verilator}Put the previous three subsections together and
a picture emerges. \SV{} will not stop you from writing a width mismatch, an
inferred latch, an incomplete sensitivity list, an unconnected port, a
multiply-driven signal or an undeclared name. The language does not consider
these to be its problem.

The industry's answer is a lint tool, run automatically, treated as a gate. The
cultural position that \VHDL{} gives to the analyser, \SV{} gives to the
linter. This is not a workaround adopted by sloppy teams. It is the mainstream
professional practice, and every serious \SV{} project you will join has it
wired into continuous integration.

Concretely, the free option is \tool{Verilator}, which is a very good linter
even if you never use it as a simulator:

\begin{shcode}
verilator --lint-only -Wall --timing -f files.f --top-module cordic_rot
\end{shcode}

\code{-Wall} turns on the checks that are off by default, including the
width-mismatch warnings (\code{WIDTHTRUNC} and \code{WIDTHEXPAND} in
\tool{Verilator}~5; older versions call both of them \code{WIDTH}),
\code{UNUSED}, \code{UNDRIVEN}, \code{CASEINCOMPLETE}, \code{LATCH} and
\code{BLKSEQ}. Commercial linters (\tool{Spyglass}, \tool{AscentLint},
\tool{Questa AutoCheck}) do the same job with more rules and a rule-set
configuration you inherit from your employer. \refsimulationbook sets
\tool{Verilator} up properly.

The other half of the answer is a written coding standard. The best public one
is the lowRISC style guide, used by the OpenTitan project; it is freely
available, it is short, and its rules are exactly the rules that stop the
failures described in this section.\index{lowRISC style guide}
\index{OpenTitan} It is worth reading in full before you write your first
thousand lines. Google's internal Verilog style guide and the ETH Zurich /
PULP platform conventions are close relatives.

\begin{ruleofthumb}
Replace the instinct ``if it compiles, the types are right'' with ``if the
linter is clean, the widths are right''. A \SV{} compile that produces no
output has told you almost nothing.
\end{ruleofthumb}

% ===========================================================================
\section{Design Units, Libraries and the Preprocessor}
\label{sec:philosophy:units}
% ===========================================================================

\index{design unit}\VHDL{} has five kinds of design unit: entity, architecture,
configuration, package and package body. Three of them are primary (entity,
package, configuration) and two are secondary (architecture, package body),
meaning they cannot be analysed until their primary exists. Each analysed unit
is deposited into a \emph{library}, which is a real directory of compiled
objects on disk, and the default library is called \code{work}. Referring to
something in another library requires you to say so twice, once to make the
library visible and once to make the name visible:

\begin{rosetta}
\begin{rleft}
library ieee;
use ieee.std_logic_1164.all;
use ieee.numeric_std.all;

library work;
use work.cordic_pkg.all;

entity cordic_rot is
  generic (ITER : positive := 14);
  port (
    clk_i  : in  std_logic;
    rst_ni : in  std_logic;
    -- ...
  );
end entity cordic_rot;

architecture rtl of cordic_rot is
  signal x_q : signed(15 downto 0);
begin
  -- ...
end architecture rtl;
\end{rleft}
\begin{rright}
import cordic_pkg::*;

module cordic_rot #(
  parameter int unsigned ITER = 14
) (
  input  logic clk_i,
  input  logic rst_ni
  // ...
);
  data_t x_q;
  // ...
endmodule
\end{rright}
\end{rosetta}

The structural differences are worth spelling out.

\textbf{One unit, not two.} A \code{module} is entity and architecture in one.
The interface and the body cannot drift apart, and there is no way to compile
two architectures of the same entity and pick between them at elaboration. The
\VHDL{} configuration has no everyday counterpart: \SV{} inherits a
\code{config}/\code{library} mechanism from Verilog-2001, but it is rarely
used, and parameterisation plus
\chOneDir{ifdef}\index{conditional compilation} do the same job. If you used
configurations to swap a behavioural model for a gate netlist, in \SV{} you do
it by choosing which files you compile.

\textbf{No library clause.} \index{compilation unit}\SV{} module names live in
a single global namespace. If two modules called \code{fifo} are compiled into
the same design, that is a collision, and the usual fix is a project-wide name
prefix, exactly as in C. What \SV{} has instead of libraries is the
\emph{compilation unit}. Everything declared outside any module lands in a
scope called \code{\$unit}. Whether \code{\$unit} covers one file or all the
files on the command line is, by the \LRM{}, tool-dependent.

\begin{gotcha}{Never declare anything at compilation-unit scope}
Because \code{\$unit} may be one file for one tool and all files for another,
a \code{typedef} or a \code{parameter} written outside a module compiles under
\tool{Verilator} and fails under \tool{VCS}, or vice versa, or, worse still,
compiles under both and means different things. Put every shared declaration in
a \kw{package} and import it. There is no \VHDL{} habit that leads you into
this, but there is a C habit that does, and \SV{} looks enough like C to
tempt you.
\end{gotcha}

\textbf{Packages exist and behave nearly as you expect.}\index{package} A
\SV{} package holds \code{typedef}s, \code{parameter}s, \code{function}s,
\code{task}s and class declarations. Import it with \code{import
cordic\_pkg::*;} for everything or \code{import cordic\_pkg::data\_t;} for one
name, and refer to individual names explicitly with
\code{cordic\_pkg::data\_t}. The scope-resolution operator \code{::} is the
counterpart of \VHDL{}'s selected name \code{work.cordic\_pkg.data\_t}. A
package must be compiled before anything that imports it, exactly as in
\VHDL{}. There is no package body: the implementations go in the package
itself.

\textbf{And there is a preprocessor.}\index{preprocessor} This is the piece
with no \VHDL{} analogue whatsoever. Verilog inherited from C a purely textual
substitution pass that runs before the language is parsed:
\chOneDir{define} creates a macro, \chOneDir{ifdef}\slash{}\chOneDir{ifndef}\slash{}
\chOneDir{else}\slash{}\chOneDir{endif} conditionally include text, and
\chOneDir{include} pastes one file into another.

\begin{svcode}[caption={Compiler directives: text substitution, not language},label={lst:philosophy:defines}]
`ifndef CORDIC_DEFS_SVH
`define CORDIC_DEFS_SVH

`define CORDIC_MAX_ITER 16
`define ASSERT(name, prop) name: assert property (prop)

`ifdef SYNTHESIS
  `define SIM_ONLY(x)
`else
  `define SIM_ONLY(x) x
`endif

`endif
\end{svcode}

What this buys you is real. Conditional compilation lets one source tree serve
simulation, synthesis, formal and emulation, with the assertion-heavy parts
excluded from the flows that cannot take them, and it lets you define something
once and use it in a hundred files without a package import. \VHDL{} has no
equivalent until the conditional analysis directives of VHDL-2019, whose tool
support is thin enough that you should not plan around it.

What it costs you is also real. Macros have no types, no scope and no
namespace; one defined anywhere is visible everywhere after that point, and
redefining it is legal. They are invisible to the elaborated design, so no
waveform or hierarchy browser will show you one. The universal conventions are
to prefix every macro name with the project name, to guard every included file
as in \cref{lst:philosophy:defines}, and to prefer a package \code{parameter}
to a \chOneDir{define} whenever a parameter will do. Use the preprocessor for
what must happen before parsing, and for nothing else.

\begin{notebox}{\chOneDir{include} is not \code{use}}
\chOneDir{include} is textual. It does not create a scope, does not check for
duplicate definitions, and can be placed inside a module body, inside a port
list, or anywhere else the resulting text would be legal. \VHDL{}'s \code{use}
is a name-visibility statement about a unit that has already been analysed. The
two look similar in a file header and have nothing in common. If a colleague
tells you to ``just include the package'', they mean import it.
\end{notebox}

% ===========================================================================
\section{The Simulation Cycle: Where Your Intuition Breaks}
\label{sec:philosophy:sched}
% ===========================================================================

This is the most important section in the chapter. Everything above is a matter
of discipline; this is a matter of semantics. If you carry your \VHDL{}
scheduling model into \SV{} unchanged, you will write designs that simulate
differently on two simulators and blame the tools.

\subsection{What VHDL does}

\index{delta cycle}You know this already, but state it precisely, because the
precision is the point of comparison.

A \VHDL{} simulation cycle has two phases in a fixed order. First the
\emph{update} phase: every signal with a transaction scheduled for the current
time $T_c$ takes its new value, and if the new value differs from the old one,
an event is posted on that signal. Second the \emph{evaluate} phase: every
process sensitive to a signal that had an event resumes, and runs until it
suspends again. Assignments a process makes to signals do not take effect
during the evaluate phase. They schedule transactions, and the earliest a
transaction can be scheduled for is $T_c$ plus one \emph{delta} -- an
infinitesimal step that does not advance simulated time.

Then the simulator computes the next time $T_n$. If $T_n = T_c$, because
somebody scheduled a delta-delayed transaction, the whole cycle repeats at the
same simulated time. That is a delta cycle. Time only advances when no
transaction remains scheduled for the current time.

Two consequences follow, and both are why \VHDL{} feels safe.

First, \emph{a signal read inside a process always returns the value it had at
the start of the current cycle}, no matter what the process assigned to it
three lines earlier. That is a rule of the language, not a convention. A
variable behaves the other way round: it updates immediately, because it is not
visible outside the process and needs no scheduling discipline.

Second, and this is the important one, \emph{a clocked \VHDL{} process is
race-free by construction}. Consider two processes, both sensitive to
\code{clk}, one assigning \code{q1 <= d} and the other \code{q2 <= q1}. Both
resume in the evaluate phase of the same cycle. Both read the pre-update values.
Both schedule for the next delta. The order in which the simulator happens to
run them cannot change the result, so a shift register is a shift register
regardless. You have never had to think about this, and that is the achievement.

\begin{notebox}{Postponed processes}
\VHDL{} also has \code{postponed process}, which runs once at the end of the
last delta cycle of a time step, when no further delta activity remains. It
exists for monitoring and reporting code that must see the settled value rather
than an intermediate delta. Keep it in mind: \SV{} has a region with exactly
this purpose, and it is even called Postponed.
\end{notebox}

\subsection{What SystemVerilog does}

\index{event regions}\SV{} does not have a two-phase cycle. It has a
\emph{time slot} containing nine regions, executed in order, with two loops
that can send execution back to an earlier region. The regions are Preponed,
Active, Inactive, NBA, Observed, Reactive, Re-Inactive, Re-NBA and Postponed.
\cref{fig:philosophy:sched} puts the two models side by side.

\begin{figure}[htbp]
  \centering
  \begin{tikzpicture}[
      font=\sffamily\scriptsize,
      bx/.style={draw=codeframe,fill=codebg,line width=0.7pt,
                 rounded corners=1.5pt,minimum width=38mm,
                 minimum height=6mm,align=center,inner sep=3pt},
      sb/.style={bx,minimum width=30mm},
      hd/.style={font=\sffamily\small\bfseries,text=codekey},
      an/.style={font=\sffamily\scriptsize\itshape,text=codenumber,
                 anchor=west,align=left,text width=32mm},
      ar/.style={-{Stealth[length=1.8mm]},line width=0.7pt,draw=codekey},
      fb/.style={-{Stealth[length=1.8mm]},line width=0.7pt,draw=accentalt,
                 dashed},
      lb/.style={font=\sffamily\tiny,text=accentalt}
    ]

    % ---------------- VHDL column ----------------
    \node[hd] at (0,1.5) {VHDL simulation cycle};
    \node[bx] (v1) at (0,0.6)  {advance time to $T_c$};
    \node[bx] (v2) at (0,-0.9) {\textbf{update phase}\\signals take new values,\\events posted};
    \node[bx] (v3) at (0,-2.6) {\textbf{evaluate phase}\\sensitive processes resume\\and run to suspension};
    \node[bx] (v4) at (0,-4.3) {assignments schedule\\transactions; compute $T_n$};
    \node[bx] (v5) at (0,-5.8) {postponed processes\\(only when $T_n > T_c$)};

    \draw[ar] (v1) -- (v2);
    \draw[ar] (v2) -- (v3);
    \draw[ar] (v3) -- (v4);
    \draw[ar] (v4) -- (v5);

    \draw[fb] (v4.west) -- (-2.75,-4.3) -- (-2.75,-0.9) -- (v2.west);
    \node[lb,rotate=90,anchor=south] at (-2.9,-2.6)
         {delta cycle: $T_n = T_c$};

    % ---------------- SV column ----------------
    \node[hd] at (6.9,1.5) {SystemVerilog time slot};
    \node[sb] (s1) at (6.9,0.6)  {Preponed};
    \node[sb] (s2) at (6.9,-0.3) {Active};
    \node[sb] (s3) at (6.9,-1.2) {Inactive};
    \node[sb] (s4) at (6.9,-2.1) {NBA};
    \node[sb] (s5) at (6.9,-3.0) {Observed};
    \node[sb] (s6) at (6.9,-3.9) {Reactive};
    \node[sb] (s7) at (6.9,-4.8) {Re-Inactive};
    \node[sb] (s8) at (6.9,-5.7) {Re-NBA};
    \node[sb] (s9) at (6.9,-6.6) {Postponed};

    \draw[ar] (s1) -- (s2);
    \draw[ar] (s2) -- (s3);
    \draw[ar] (s3) -- (s4);
    \draw[ar] (s4) -- (s5);
    \draw[ar] (s5) -- (s6);
    \draw[ar] (s6) -- (s7);
    \draw[ar] (s7) -- (s8);
    \draw[ar] (s8) -- (s9);

    \node[an] at (8.6,0.6)  {sample values for\\assertions and\\clocking inputs};
    \node[an] at (8.6,-0.3) {blocking assignments,\\continuous assignments,\\RHS of \texttt{<=}};
    \node[an] at (8.6,-1.2) {\texttt{\#0} delays};
    \node[an] at (8.6,-2.1) {left sides of \texttt{<=}\\are updated};
    \node[an] at (8.6,-3.0) {concurrent assertions\\evaluated};
    \node[an] at (8.6,-3.9) {\texttt{program} blocks,\\assertion actions};
    \node[an] at (8.6,-5.7) {clocking-block\\outputs driven};
    \node[an] at (8.6,-6.6) {\texttt{\$monitor},\\\texttt{\$strobe}};

    \draw[fb] (s4.west) -- (5.05,-2.1) -- (5.05,-0.3) -- (s2.west);
    \node[lb,rotate=90,anchor=south] at (4.9,-1.2) {iterate};
    \draw[fb] (s8.west) -- (4.35,-5.7) -- (4.35,-0.3) -- (s2.west);
    \node[lb,rotate=90,anchor=south] at (4.2,-3.5) {re-enter};
  \end{tikzpicture}
  \caption{The \VHDL{} simulation cycle and the \SV{} time slot. \VHDL{}
    separates update from evaluate once, and repeats the pair in delta cycles.
    \SV{} splits the same job across nine regions and iterates between them
    until no events remain.}
  \label{fig:philosophy:sched}
\end{figure}

Region by region, in the terms you already have:

\begin{description}[leftmargin=1.6em,style=nextline,font=\normalfont\ttfamily]
\item[Preponed] Nothing changes value here. Values are \emph{sampled} for
  concurrent assertions and for clocking-block inputs. This is what
  \code{\#1step} means: sample as the value was immediately before this time
  slot.
\item[Active] Where ordinary procedural code runs. Blocking assignments take
  effect, continuous assignments are evaluated, \code{\$display} prints, and
  the \emph{right-hand sides} of non-blocking assignments are evaluated and
  queued. Processes here execute in an order the \LRM{} deliberately leaves
  unspecified.
\item[Inactive] Processes that executed an explicit \code{\#0} delay.
\item[NBA] The non-blocking assignment queue is applied: the left-hand sides
  now take the values that were computed in Active. This is \VHDL{}'s update
  phase.
\item[Observed] Concurrent assertion properties are evaluated, using the values
  sampled in Preponed.
\item[Reactive] Testbench code in a \code{program} block, and the action blocks
  of assertions.
\item[Re-Inactive, Re-NBA] The same as Inactive and NBA, for processes in the
  reactive set. Clocking-block outputs are driven in Re-NBA.
\item[Postponed] \code{\$monitor} and \code{\$strobe}. Read-only: nothing may
  change a value here.
\end{description}

The loops matter as much as the regions. After NBA, if anything that just
happened created new active events, execution returns to Active and the
Active/Inactive/NBA group repeats until all three are empty. This iteration is
the exact counterpart of a chain of \VHDL{} delta cycles. There is a second
loop from the reactive group back to Active, so a testbench that drives a
signal can wake design code inside the same time slot.

\begin{notebox}{The regions the \LRM{} lists that this book ignores}
IEEE 1800-2017 Figure 4-1 actually shows more than nine regions: Pre-NBA,
Post-NBA, Post-Observed, Pre-Re-NBA and Post-Re-NBA exist as well. They are
reserved for PLI and VPI callbacks, so that a C model attached to the simulator
can run at a defined point relative to the language's own events. No \SV{} you
write ever executes in them, and they are omitted from
\cref{fig:philosophy:sched} for clarity.
\end{notebox}

\subsection{Blocking, non-blocking, and why Verilog needed them}

\index{blocking assignment}\index{non-blocking assignment}Now the pieces join
up. \SV{} has two procedural assignment operators:

\begin{itemize}
\item \code{=}, the \emph{blocking} assignment. It evaluates its right-hand
  side and updates its target immediately, in the Active region, before the
  next statement in the process runs. It behaves like a \VHDL{}
  \kw{variable}.
\item \code{<=}, the \emph{non-blocking} assignment. It evaluates its
  right-hand side in the Active region, queues the result, and updates its
  target later, in the NBA region. It behaves like a \VHDL{} \kw{signal} with a
  delta delay.
\end{itemize}

That is the whole mapping, and it is exact enough to rely on:

\begin{center}
\begin{tabular}{@{}ll@{}}
\toprule
\textbf{\VHDL{}} & \textbf{\SV{}} \\
\midrule
\code{variable v : ...;\quad v := expr;} & \code{logic v;\quad v = expr;} \\
\code{signal s : ...;\quad s <= expr;}   & \code{logic s;\quad s <= expr;} \\
\bottomrule
\end{tabular}
\end{center}

Here is the historical point that makes the whole design comprehensible.
\VHDL{} never needed to invent the non-blocking assignment, because the signal
\emph{already} had a delta-delayed update built into its definition: there is
only one kind of assignment to a signal, and it is always deferred. Verilog
started from the opposite place, with an immediate procedural assignment and
\code{always} blocks scheduled in an unspecified order, so writing a shift
register was genuinely ambiguous. The non-blocking assignment was added to
Verilog-XL, and standardised in IEEE 1364-1995, precisely to recover the
discipline \VHDL{} had from birth: evaluate everything, then update
everything.

Look at the same register in both languages, and read the right-hand pane as
``the \code{<=} is doing what your \VHDL{} signal did automatically'':

\begin{rosetta}
\begin{rleft}
regs : process (clk_i, rst_ni) is
begin
  if rst_ni = '0' then
    q <= (others => '0');
  elsif rising_edge(clk_i) then
    q <= d;
  end if;
end process regs;
\end{rleft}
\begin{rright}
always_ff @(posedge clk_i or
            negedge rst_ni) begin
  if (!rst_ni) q <= '0;
  else         q <= d;
end
\end{rright}
\end{rosetta}

Eight lines against five, and the five say the same thing. The \code{'0} is
\SV{}'s unsized fill literal, the counterpart of \code{(others => '0')}: it
fills the target with zeros whatever its width. Note that the \SV{} sensitivity
list is \code{negedge rst\_ni}, not the plain \code{rst\_ni} of the \VHDL{}
process. \SV{} requires every term of an \code{always\_ff} event control to be
edge-qualified, and it needs the polarity to know which edge starts the
asynchronous reset. \VHDL{} works out the same thing from the \code{if}
statement.

\subsection{The classic race}

Now the failure that this machinery exists to prevent.
\cref{lst:philosophy:race} is a two-stage shift register written with blocking
assignments.

\begin{svcode}[caption={A race. This design has two legal outcomes.},label={lst:philosophy:race}]
module shift_bad (input logic clk, input logic d, output logic q2);
  logic q1;
  always_ff @(posedge clk) q1 = d;    // blocking: wrong
  always_ff @(posedge clk) q2 = q1;   // blocking: wrong
endmodule
\end{svcode}

Both processes wake in the Active region of the same time slot. The \LRM{} does
not say which runs first. If the first runs first, \code{q1} takes the value of
\code{d} immediately, then the second reads the \emph{new} \code{q1} and
\code{q2} follows \code{d} in a single cycle: one flop, not two. If the second
runs first, \code{q2} takes the old \code{q1} and you get the two-stage shift
register you intended. Both are legal. Two simulators may disagree; one
simulator may disagree with itself after you add an unrelated module and change
its internal scheduling order. Synthesis, meanwhile, will happily produce two
flip-flops, so the gate-level netlist need not match the RTL simulation.

\begin{gotcha}{The failure mode you have never had to think about}
In \VHDL{} you could write the two processes in either order, in either file,
and get the same answer every time, because signal semantics forbid the race.
In \SV{} the ordering of \code{always} blocks within the Active region is
explicitly unspecified, so the race is real and the language will not warn you.
Change one character in \cref{lst:philosophy:race}:

\begin{svcode}
  always_ff @(posedge clk) q1 <= d;
  always_ff @(posedge clk) q2 <= q1;
\end{svcode}

Now both right-hand sides are evaluated in Active, using the old values, and
both targets are updated in NBA. The order of the two blocks is irrelevant.
This is why the rule exists, and it is why \tool{Verilator}'s \code{BLKSEQ}
warning and every commercial lint rule set flag a blocking assignment inside a
clocked block as an error, not a style issue.
\end{gotcha}

The rule that follows is the first thing anyone will tell you about \SV{}, and
now you know why it is true rather than merely being told it.

\begin{ruleofthumb}
Non-blocking (\code{<=}) for everything a clock edge writes. Blocking
(\code{=}) for everything combinational. Never mix the two in one
\code{always} block, and never write the same variable from two blocks.
\end{ruleofthumb}

The reasoning for the combinational half is the mirror image. Inside an
\code{always\_comb}, a blocking assignment behaves like a \VHDL{} variable in a
combinational process: intermediate results are visible to later statements in
the same execution, which is exactly what you want when you are describing a
chain of logic. Using \code{<=} there would push the update into the NBA
region, so a later statement in the same block would read the stale value, and
the block would then have to be re-executed, producing extra delta activity,
slower simulation and, in some cases, a mismatch with synthesis.

\subsection{Testbenches race too}

\index{clocking block}The design-side race has a clean rule. The
testbench-side race does not, and it needs language support.

Consider a testbench driving a synchronous input:

\begin{svcode}
initial begin
  @(posedge clk);
  data = 8'h5A;    // Active region, same time slot as the DUT's clock edge
end
\end{svcode}

The \DUT samples \code{data} on the same edge, and the testbench writes it in
the Active region of that same time slot, so whether the \DUT sees the old
value or the new one depends on which process the simulator happens to run
first. The ``use non-blocking'' rule does not fully save you here, because the
testbench also wants to \emph{read} \DUT outputs at the clock edge, and
reading has the same ambiguity. Verilog testbench writers solved this by hand
for twenty years, inserting small delays after the clock edge. \SV{} solved it
in the language, with the \emph{clocking block}: a declaration that names a
clock and states, once, when its signals are sampled and when they are
driven.

\begin{svcode}[caption={A clocking block removes the testbench race by construction},label={lst:philosophy:clkblk}]
module tb;
  logic       clk = 1'b0;
  logic       vld, ready;
  logic [7:0] data;

  always #5ns clk = ~clk;

  clocking cb @(posedge clk);
    default input #1step output #1ns;
    output vld, data;
    input  ready;
  endclocking

  initial begin
    @(cb);
    cb.vld  <= 1'b1;         // driven in Re-NBA, 1 ns after the edge
    cb.data <= 8'h5A;
    @(cb);
    if (cb.ready) $display("accepted");  // sampled 1step before the edge
  end
endmodule
\end{svcode}

\code{input \#1step} means every clocking-block input is sampled in the
Preponed region, before any design code in that time slot has run, so it can
never see a value the \DUT produced at this very edge. \code{output \#1ns}
means every drive happens 1\,ns after the edge, safely inside the cycle. The
race is gone by construction, and it stays gone when somebody adds another
\code{always} block to the design six months later. \refverificationbook builds
testbenches on this foundation.

\begin{notebox}{Whatever happened to \code{program} blocks?}
\index{program block}\SV{} 3.1 introduced the \kw{program} block, borrowed from
Vera: a container whose code runs in the Reactive region, which is after all
design code in the same time slot, so testbench statements could never race
with the design. It solved a real problem, but it also brought restrictions
(no \code{always} blocks inside, particular termination semantics) and it was
largely superseded by clocking blocks, which fix the same races without
constraining the testbench structure. UVM does not use \code{program} blocks.
You will meet them in older code and in textbooks; you do not need to write
them.
\end{notebox}

\begin{gotcha}{\code{after} is not \code{\#}}
\VHDL{}'s \code{s <= v after 5 ns} schedules a transaction, with inertial
delay: a pulse shorter than 5\,ns is swallowed. \SV{}'s \code{\#} is not a
scheduling annotation on an assignment, it is a statement delay, and its two
placements mean different things. \code{a = \#5 b} evaluates \code{b} now,
suspends the process for 5\,ns and then assigns; the process is blocked for
those 5\,ns. \code{a <= \#5 b} evaluates \code{b} now and schedules the update
for 5\,ns later without blocking. Neither is inertial. And in synthesisable
RTL you should write no \code{\#} delays at all -- the closest \SV{} gets to
\VHDL{}'s inertial behaviour is on gate primitives and continuous assignments,
which is not where you will be working.
\end{gotcha}

% ===========================================================================
\section{Intent, Declared and Checked}
\label{sec:philosophy:intent}
% ===========================================================================

\SV{} gave away the type safety that \VHDL{} has. It got something back that
\VHDL{} does not have, and the trade is more even than it first looks. In
\VHDL{} a process is a process: whether it describes a flip-flop, combinational
logic or a latch is something the reader, the simulator and the synthesiser
each work out from the shape of the code inside it. There is no place to write
down which of the three you meant, so no tool can tell you that you wrote
something else. Every \VHDL{} designer has produced a latch by forgetting an
\code{else} in a combinational process, and the analyser said nothing, because
nothing was wrong with the \emph{types}.

\SV{}-2005 added three specialised process keywords whose only purpose is to
state intent so that a tool can check it.\index{always\_ff}
\index{always\_comb}\index{always\_latch}

\begin{itemize}
\item \code{always\_ff} declares sequential logic. Its event control must
  contain only edge-qualified terms; the variables it writes may not be written
  anywhere else; and tools check that what it describes is in fact a
  flip-flop.
\item \code{always\_comb} declares combinational logic. Its sensitivity list is
  inferred from everything the block reads, \emph{including inside functions it
  calls}; the variables it writes may not be written anywhere else; and a tool
  will warn if the code as written infers a latch. It also executes once at the
  end of time zero, whether or not any input changed, so the outputs are
  consistent from the start of the simulation.
\item \code{always\_latch} declares a latch on purpose, so that a real latch
  does not look like a mistake.
\end{itemize}

\begin{rosetta}
\begin{rleft}
-- intent is not stated; a missing
-- else silently infers a latch
comb : process (all) is
begin
  if sel = '1' then
    y <= a;
  else
    y <= b;
  end if;
end process comb;
\end{rleft}
\begin{rright}
// intent is stated and checked
always_comb begin
  if (sel) y = a;
  else     y = b;
end
\end{rright}
\end{rosetta}

The \VHDL{}-2008 \code{process (all)} form removes the incomplete sensitivity
list problem, which is the other half of the same class of bug, so that part is
a draw. What \VHDL{} still has no way to express is the sentence ``this block
is combinational, please complain if it is not''.

Argue it this way. \VHDL{} applies static checking to \emph{data}: what a value
is, how wide it is, what may be assigned to what. \SV{} applies static checking
to \emph{structure}: what kind of hardware a block of code describes, and
whether anything else drives the same variable. Neither language checks what
the other checks. A \VHDL{} designer arriving in \SV{} should use
\code{always\_ff} and \code{always\_comb} everywhere with the same reflex that
made them declare a signal as \code{unsigned(11 downto 0)} rather than
\code{std\_logic\_vector}: it is free, and it converts a class of silent bugs
into a message.

\begin{tipbox}{Naming, and why the industry cares about it}
\SV{} has no reserved suffix conventions, so a project's naming rules do a job
that the language does not. Two notes before the list. These rules apply to
design code you intend to keep; the small throwaway examples in this book ---
a counter, a FIFO, a shift register --- drop the suffixes so that the line
under discussion fits in a narrow column, and the CORDIC sources of
\refexamplebook follow the rules in full. And the rules themselves matter less
than picking one set and holding to it. The set below comes from the lowRISC guide and is
close to universal in modern \SV{}:

\begin{itemize}
\item \code{\_i}, \code{\_o}, \code{\_io} on module ports for input, output and
  bidirectional. The direction is then visible at every use, not only at the
  declaration.
\item \code{\_n} for active-low, always as the last element before the
  direction: a reset input is \code{rst\_ni}.
\item \code{\_d} and \code{\_q} for the input and the output of a register.
  \code{cnt\_d} is what will be stored, \code{cnt\_q} is what is stored. An
  \code{always\_ff} then contains nothing but \code{cnt\_q <= cnt\_d;} and all
  the logic lives in an \code{always\_comb} that computes \code{cnt\_d}.
\item Lower case with underscores for signals and instances, upper case for
  \code{parameter}s and macros, \code{\_t} for a \code{typedef}, \code{\_e} for
  an \code{enum}, \code{\_pkg} for a package.
\end{itemize}

The \code{\_d}/\code{\_q} convention deserves a specific defence: it makes the
register boundary visible in the source. In \VHDL{} you can usually see it
because the clocked process is a separate lump of text; in a language where
\code{always\_ff @(posedge clk\_i) x <= y;} is one line, you need the names to
carry the information.
\end{tipbox}

\begin{tipbox}{One clock, one reset, one \code{always\_ff}}
Give every \code{always\_ff} exactly one clock and at most one asynchronous
reset. Do not put logic for two clock domains in one block, do not gate a clock
inside the event control, and do not use a signal as both a clock and data.
Clock-domain crossings get their own dedicated modules with their own lint and
CDC checks. This is the same discipline you already apply in \VHDL{}; \SV{}
simply makes it easier to violate, because the terse syntax invites you to put
more into one block.
\end{tipbox}

\subsection{Brevity is a feature}

The register comparison above was eight lines against five. That ratio holds
across most RTL, and it is not only about typing speed. Verbosity has a real
cost: a design that fits on one screen can be checked by eye, and a design that
takes three screens cannot. When the same finite-state machine is 120 lines
instead of 300, more of it is visible at once and structural mistakes stand
out.

The savings come from several places at once: no separate entity, no
\code{architecture} wrapper, no declarative region separate from the body, no
explicit type conversions, \code{begin}/\code{end} omitted for a single
statement, ANSI port lists with types inline, and the \code{'0} fill literal.
Nothing on that list is free, since every item is an explicitness you have
given up, but the aggregate effect on readability is positive and it is one of
the honest reasons the industry moved.

% ===========================================================================
\section{Verification: the Bigger Break}
\label{sec:philosophy:verif}
% ===========================================================================

Everything so far has been about writing hardware. The larger difference
between the two ecosystems is about proving that the hardware is right, and
here the gap is not a matter of syntax at all. It is a different theory of what
a testbench \emph{is}.

\subsection{The VHDL testbench you already write}

\index{testbench!VHDL}In \VHDL{}, a testbench is another design unit. It is an
entity with an empty port list and an architecture that instantiates the device
under test and drives it from one or more processes. The stimulus is
procedural, directed and written by hand: assert reset, wait, apply a value,
wait for a clock edge, check the output with \code{assert ... report ...
severity ...}. Reference data comes from a file read with \code{TEXTIO} or from
a function written in the testbench.

\begin{vhdlcode}[caption={The shape of a directed VHDL testbench},label={lst:philosophy:vhdltb}]
stim : process is
begin
  rst_ni <= '0';
  wait for 40 ns;
  rst_ni <= '1';
  wait until rising_edge(clk_i);

  req_x_i     <= to_signed(16#2000#, 16);
  req_y_i     <= (others => '0');
  req_valid_i <= '1';
  wait until rising_edge(clk_i) and req_ready_o = '1';
  req_valid_i <= '0';

  wait until rising_edge(clk_i) and rsp_valid_o = '1';
  assert abs(to_integer(rsp_x_o) - 11585) < 4
    report "cos result out of tolerance" severity error;

  std.env.stop;
end process stim;
\end{vhdlcode}

There is nothing wrong with this. It is readable, deterministic and easy to
debug, and for a small block with a handful of interesting cases it is the
right answer. Its limitation is arithmetical rather than aesthetic: you verify
exactly the cases you thought of, and the number of cases you can hand-write
does not grow with the size of the design, while the number of cases the design
can get wrong does. The language gives you nothing beyond this: \VHDL{} has no
classes, no inheritance, no constrained randomisation and no functional
coverage constructs, only \code{ieee.math\_real.uniform}, which is a random
number generator rather than a stimulus generator.

\subsection{Verification as a software problem}

\index{constrained random}\index{functional coverage}The \SV{} verification
half of the language starts from a different premise: verifying a large design
is a software engineering problem, so give the verification engineer software
engineering tools. Four ideas follow, and each one changes what a testbench
looks like. The constructs behind them are \refverificationbook; what follows here
is only the argument for why they exist.

\textbf{Transactions instead of pin wiggles.} You stop writing values onto
individual signals and start describing a \emph{transaction}: a read of four
bytes from address \code{0x40}, a rotation of a vector by an angle. A separate
component, the driver, knows how to turn a transaction into pin activity for a
particular protocol. The test then says what should happen, not how the wires
move, and the same test survives a change of bus protocol.

\textbf{Constrained random instead of directed vectors.} Instead of enumerating
inputs, you describe the \emph{legal input space} and let the tool draw from
it. This is the deepest change of habit in the whole book.

\textbf{Functional coverage as the exit criterion.} If the stimulus is random,
``all my tests passed'' means nothing on its own, because you no longer know
what was tried. So you write down, separately and explicitly, the situations
that must occur, and the simulator counts them. A run is worth something when
the coverage model says the interesting cases happened.

\textbf{Object-oriented reuse.} Drivers, monitors, scoreboards and sequences
are classes, and a new project extends them rather than copying them. This is
what UVM is: not an extension of the language but a class library written in
\SV{}, which \refverificationbook takes apart.

\cref{lst:philosophy:crv} is a minimal version of the first three ideas.

\begin{svcode}[caption={Constraints describe the legal space; coverage measures what happened},label={lst:philosophy:crv}]
class rot_req;
  rand bit signed [15:0] x, y;
  rand bit signed [15:0] theta;
  rand bit               small_angle;

  // the legal input space, written once
  constraint c_range {
    x inside {[-16384:16383]};
    y inside {[-16384:16383]};
    small_angle -> theta inside {[-1024:1023]};
  }
endclass

covergroup cg_req with function sample(bit signed [15:0] theta);
  cp_theta : coverpoint theta {
    bins neg_large = {[-32768:-8193]};
    bins neg_small = {[-8192:-1]};
    bins zero      = {0};
    bins pos_small = {[1:8192]};
    bins pos_large = {[8193:32767]};
  }
endgroup
\end{svcode}

The constraint block is declarative. It does not say what to do; it says what
is true of every legal input. Calling \code{req.randomize()} asks a constraint
solver for a value satisfying all of it, and calling it ten thousand times with
different seeds explores parts of the space you would never have enumerated by
hand. The \code{covergroup} then records what actually came out, so you can
tell the difference between ``ten thousand runs'' and ``ten thousand runs that
all did the same thing''.

\subsection{The metric-driven loop}

The two flows differ in shape, not only in tooling.
\cref{fig:philosophy:flows} draws them.

\begin{figure}[htbp]
  \centering
  \begin{tikzpicture}[
      font=\sffamily\scriptsize,
      bx/.style={draw=codeframe,fill=codebg,line width=0.7pt,
                 rounded corners=1.5pt,minimum width=34mm,
                 minimum height=7mm,align=center,inner sep=3pt},
      dm/.style={draw=codeframe,fill=tipbg,line width=0.7pt,
                 diamond,aspect=2.1,align=center,inner sep=1pt,
                 minimum width=30mm},
      dn/.style={bx,fill=tipbg},
      hd/.style={font=\sffamily\small\bfseries,text=codekey},
      ar/.style={-{Stealth[length=1.8mm]},line width=0.7pt,draw=codekey},
      lb/.style={font=\sffamily\tiny,text=accentalt}
    ]
    \node[hd] at (0,1.2) {Directed};
    \node[bx] (d1) at (0,0.3)  {list the cases you\\can think of};
    \node[bx] (d2) at (0,-1.3) {write one test each};
    \node[bx] (d3) at (0,-2.9) {run them all};
    \node[dn] (d4) at (0,-4.5) {they pass};
    \node[bx] (d5) at (0,-6.1) {sign off\\(what was missed\\is unknown)};
    \draw[ar] (d1) -- (d2);
    \draw[ar] (d2) -- (d3);
    \draw[ar] (d3) -- (d4);
    \draw[ar] (d4) -- (d5);

    \node[hd] at (7.5,1.2) {Coverage-driven};
    \node[bx] (c1) at (7.5,0.3)  {describe the legal\\input space as\\constraints};
    \node[bx] (c2) at (7.5,-1.3) {write the coverage\\model: what must\\happen};
    \node[bx] (c3) at (7.5,-2.9) {run many seeds;\\a scoreboard checks\\every result};
    \node[dm] (c4) at (7.5,-4.6) {coverage\\complete?};
    \node[bx] (c5) at (7.5,-6.4) {sign off against\\a measured metric};
    \draw[ar] (c1) -- (c2);
    \draw[ar] (c2) -- (c3);
    \draw[ar] (c3) -- (c4);
    \draw[ar] (c4) -- node[lb,right,pos=0.4] {yes} (c5);
    \draw[ar] (c4.west) -- (4.9,-4.6) -- (4.9,0.3) -- (c1.west);
    \node[lb,rotate=90,anchor=south] at (4.75,-2.1)
         {no: add constraints, add seeds};
  \end{tikzpicture}
  \caption{Directed verification stops when the tests you wrote pass.
    Coverage-driven verification stops when a measurement says the input space
    has been exercised. Only the second one has a closed loop.}
  \label{fig:philosophy:flows}
\end{figure}

The mental shift is from ``I will write a test for every case I can think of''
to ``I will describe the legal input space and let the tool find the cases''.
The first is bounded by your imagination. The second is bounded by your
coverage model, which is a written artefact that a colleague can review and
that a project manager can track. When a bug escapes a coverage-driven flow,
the post-mortem question is precise and answerable: which coverage point would
have caught it, and why was it not in the model?

\begin{gotcha}{Random stimulus without a scoreboard proves nothing}
A directed test carries its expected result in the same file. A random test
cannot. If you randomise the inputs, something else has to know what the right
answer is: a reference model, a golden function, a protocol checker, an
assertion. Newcomers write the constraints first because they are the fun part,
run a million cycles, see no errors and conclude that the design works. The
design was never checked. Build the checker before the generator, always.
\end{gotcha}

\subsection{Be fair to VHDL}

None of the above is impossible in \VHDL{}, and it would be dishonest to
suggest otherwise. \index{OSVVM}\index{UVVM}OSVVM and UVVM are mature,
maintained, free \VHDL{} verification libraries. OSVVM's \code{RandomPkg} gives
randomisation with weighted distributions and \code{CoveragePkg} gives
functional coverage with bins and goals; both libraries add scoreboards and
transaction-level interfaces, and both are used in production, particularly in
European aerospace and industrial companies.

\begin{vhdlcode}[caption={OSVVM: constrained random and coverage in VHDL (sketch)},label={lst:philosophy:osvvm}]
library osvvm;
use osvvm.RandomPkg.all;
use osvvm.CoveragePkg.all;
-- ...
stim : process is
  variable rv  : RandomPType;
  variable cov : CoveragePType;
begin
  cov.AddBins(GenBin(-32768, 32767, 5));   -- five bins over the range
  while not cov.IsCovered loop
    theta := rv.RandInt(-32768, 32767);
    drive_transaction(theta);              -- a procedure, not pin wiggles
    cov.ICover(theta);
  end loop;
  cov.WriteBin;
  std.env.stop;
end process stim;
\end{vhdlcode}

So the honest statement of the difference is this. It is not that \VHDL{}
cannot do coverage-driven verification. It is that in \SV{} these constructs
are in the language, in the \LRM{}, in every commercial simulator, in the
university curriculum and in the job advertisements, whereas in \VHDL{} they
are a library you must choose, install and persuade your team to adopt. The
difference is one of ecosystem and industrial default rather than of
possibility. That distinction matters when you are deciding what to learn,
because the ecosystem is what you will be hired into.

\begin{notebox}{Assertions on both sides}
\index{assertion}\VHDL{}'s \code{assert ... report ... severity} is an
immediate, procedural check: it tests a condition at the moment the statement
executes. \SV{} has the same thing (\code{assert (cond) else \$error(...);})
and adds \emph{concurrent} assertions, which describe behaviour over time
using a temporal sequence language: \code{assert property (@(posedge clk\_i)
req\_valid\_i \&\& !req\_ready\_o |=> \$stable(req\_x\_i));}. Those are
evaluated in the Observed region on values sampled in Preponed, which is
exactly why the scheduler section came first. The sequence language is a
chapter of its own: \refverificationbook. \VHDL{} can express the same properties
with PSL, which \VHDL{}-2008 integrated into the language, but PSL support
varies between tools and the construct is much less commonly taught.
\end{notebox}

% ===========================================================================
\section{The Contrasts in One Table}
\label{sec:philosophy:summary}
% ===========================================================================

\cref{tab:philosophy:summary} collects the argument. Read it as a summary, not
as a reference: every row is a topic that a later chapter treats properly.

\begin{table}[htbp]
  \centering
  \caption{The philosophical contrasts, summarised.}
  \label{tab:philosophy:summary}
  \small
  \begin{tabularx}{\linewidth}{@{}lXX@{}}
    \toprule
    & \textbf{\VHDL{}} & \textbf{\SV{}} \\
    \midrule
    Typing
      & Strong and static. Nearly everything is checked at analysis or
        elaboration.
      & Weak by default. Widths and signedness are adjusted silently. \\
    Conversions
      & Never implicit. \code{to\_integer}, \code{unsigned()},
        \code{resize} are compulsory.
      & Almost always implicit. Casts exist and are a matter of discipline. \\
    Where errors are caught
      & The analyser.
      & The linter, the coding standard and simulation. \\
    Design units
      & Entity, architecture, configuration, package, package body.
      & One \code{module}. Packages exist; configurations do not, in
        practice. \\
    Namespaces
      & Libraries with \code{work} as default; \code{library}/\code{use}.
      & One global namespace plus packages and \code{::}; the compilation
        unit is best left empty. \\
    Preprocessor
      & None before VHDL-2019, whose support is thin.
      & Full textual preprocessor: \chOneDir{define}, \chOneDir{ifdef},
        \chOneDir{include}. \\
    Concurrency model
      & Signals with delta-delayed update; two-phase cycle.
      & Nine event regions per time slot, with iteration. \\
    Deferred update
      & Built into the \kw{signal}; there is only one kind of signal
        assignment.
      & Chosen per statement: \code{<=} defers, \code{=} does not. \\
    Race handling
      & Race-free between processes by construction.
      & Race-free only if you follow the assignment rules and use clocking
        blocks. \\
    Declared intent
      & None. A process is a process.
      & \code{always\_ff}, \code{always\_comb}, \code{always\_latch}, checked
        by the tool. \\
    Genericity
      & Generics, generate, unconstrained arrays, VHDL-2008 generic
        packages.
      & Parameters, \code{generate}, parameterised interfaces and classes. \\
    Verification constructs
      & In the language: none beyond \code{assert}/\code{report}. In
        libraries: OSVVM, UVVM.
      & In the language: classes, \code{rand}, constraints, \code{covergroup},
        SVA, mailboxes. \\
    Standard
      & IEEE 1076; current 2008, with 2019 poorly supported.
      & IEEE 1800; current 2017/2023. Absorbed IEEE 1364 in 2009. \\
    Free tools
      & Mature free simulation of the whole language. Synthesis via
        \tool{Yosys} plugins.
      & \tool{Verilator} and \tool{Yosys}: fast, but a language subset
        (\refsimulationbook). \\
    Industrial default
      & European aerospace, defence, FPGA houses, academia.
      & ASIC design and verification worldwide; UVM is the de facto
        methodology. \\
    \bottomrule
  \end{tabularx}
\end{table}

One row deserves an extra sentence. ``Free tools'' is not a footnote for this
book: \tool{Verilator} is the reason you can do serious \SV{} work on a
laptop, and \refsimulationbook is built on it. Be aware
of its boundaries, though. \tool{Verilator} compiles a language subset into
C++; it is not a general-purpose event simulator, it is two-state, and it
implements neither covergroups nor the constraint solver, so UVM is out of
reach. \refsimulationbook states this precisely and lists what works where. The row is picked up again in
\cref{sec:philosophy:vhdlwins}, where the free-tool position is one of the
places \VHDL{} is ahead.

% ===========================================================================
\section{What Actually Changes in Your Habits}
\label{sec:philosophy:habits}
% ===========================================================================

Abstractions are easier to remember when they come with a list of
behaviours. Here is what to change on day one.

\subsection{Stop doing}

\begin{itemize}
\item \textbf{Stop trusting a clean compile.} No messages means the file
  parsed. It does not mean the widths agree, the ports are connected, the block
  is combinational or the name existed. A file is not finished until the linter
  is clean too.
\item \textbf{Stop assuming a name is declared.} Until
  \chOneDir{default\_nettype none} is in the file, an undeclared name is a
  one-bit wire, not an error.
\item \textbf{Stop expecting a range check.} There is no \code{natural}, no
  \code{range 0 to 9} and no subtype constraint, so nothing complains when a
  value leaves the range you had in mind. If it must be checked, assert it.
\item \textbf{Stop reaching for two-state types because they look tidy.}
  \kw{int} and \kw{bit} in RTL hide exactly the reset bugs that \code{'U'}
  propagation used to reveal.
\end{itemize}

\subsection{Start doing}

\begin{itemize}
\item \textbf{Start with the directive pair.} Put
  \chOneDir{default\_nettype none} at the top of every file, and
  \chOneDir{default\_nettype wire} at the bottom of it.
\item \textbf{Start using \code{always\_ff} and \code{always\_comb}
  exclusively.} Plain \code{always} in new RTL is a smell, and
  \code{always @(*)} is a worse one.
\item \textbf{Start naming with suffixes:} \code{\_i}, \code{\_o},
  \code{\_ni}, \code{\_d}, \code{\_q}, \code{\_t}, \code{\_e}. It costs nothing
  and it makes a code review possible.
\item \textbf{Start running the linter from the first hour.} Put
  \code{verilator --lint-only -Wall} in your \file{Makefile} before the second
  module, not after the first bug; \refsimulationbook shows the full setup.
\item \textbf{Start reading the \LRM{}, clause 4.} Ten pages on scheduling
  semantics will save you a week over a project.
\item \textbf{Start writing the checker before the stimulus.}
  \refverificationbook builds testbenches in that order for a reason.
\end{itemize}

\subsection{Be paranoid about, for the first month}

\begin{itemize}
\item \textbf{Every width:} every assignment, port connection, concatenation
  and constant. Read the lint report before you open the waveform.
\item \textbf{Every \code{=} inside a clocked block, and every \code{<=} inside
  a combinational one.} Both are legal and both are wrong.
\item \textbf{Every signed operation.} If any operand of an expression is
  unsigned, the whole expression is evaluated unsigned, so one unsigned
  constant changes the result. \cref{ch:rosetta} treats this properly.
\item \textbf{Every port left unconnected.} \code{.*} and omitted ports are
  legal and quiet.
\item \textbf{Every testbench signal driven at a clock edge.} If it is not
  inside a clocking block, assume it races until you have proved otherwise.
\item \textbf{Every \code{case} without a \code{default}.} There is no
  \code{when others} obligation in \SV{}: write the \code{default}, or use
  \code{unique case} and let the tool check it.
\end{itemize}

The list is not a list of hard things. It is a list of things the language does
not check for you, and almost every one is caught by a lint rule that somebody
else has already written, in a configuration file you will inherit on your
first day of work. The point of this chapter is to explain why those rules
exist, so that when one fires you know what it is protecting you from rather
than looking for the pragma that switches it off.

Which brings the argument to its other half. Everything up to here has been
about what changes, and mostly about what you gain, when you move to \SV{}.
The two sections that close the chapter are about what you leave behind, and
about how to decide, for a particular project, whether leaving it behind is
the right trade. The honest answer to the first question is: more than a book
with this title normally admits.

% ===========================================================================
\section{What \VHDL{} has that \SV{} does not}
\label{sec:philosophy:vhdlwins}
\label{sec:beyond:vhdlwins}

Almost everything above is a case for \SV{}, or at least an explanation of it.
This section is the case against, and it is not a short list. Read it before
you argue that one language is simply better than the other, and read it again
before you propose a language change to a team that has a working \VHDL{}
flow.

\subsection{A type system that catches errors at compile time}
\index{type system}

This is the big one, and \cref{sec:philosophy:typing} has already made the
core of the argument: widths and signedness that \VHDL{} checks at analysis
time, \SV{} adjusts silently, and the compensation is a lint tool rather than
a compiler. Two consequences of that were not stated there.

The first is about your own types. \VHDL{} lets you define a scalar type that
is genuinely new, an address type that is not interchangeable with a data type
of the same width, and the compiler enforces the distinction. \SV{} has
\code{typedef}, but a typedef is an alias, not a new type: assigning an
\code{addr\_t} to a \code{data\_t} of the same width is legal and silent. All
the discipline a \VHDL{} designer gets from naming a type has to come from
review instead.

The second is about arithmetic in the large. \VHDL{}'s \code{signed} and
\code{unsigned} are distinct types and mixing them requires a conversion you
have to write down. In \SV{}, signedness is a property of an expression,
computed by rules that involve every operand, so one unsigned intermediate
three terms away turns an arithmetic right shift into a logical one.
\cref{ch:rosetta} spends pages on those rules for a reason, and the reason is
that they are the price of this section's first paragraph.

\begin{ruleofthumb}
The single biggest practical loss in moving to \SV{} is compile-time width
checking. Compensate with lint. Run a lint tool on every commit and treat
width warnings as errors, because your compiler will not.
\end{ruleofthumb}

\subsection{Physical types with units}
\index{physical type}

\VHDL{} has physical types: quantities with units and a compiler that checks
them. \code{time} is the built-in one, and you can declare your own:

\begin{vhdlcode}
type capacitance is range 0 to 1e9
  units
    ff;  pf = 1000 ff;  nf = 1000 pf;
  end units;

constant t_setup : time := 250 ps;
wait for t_setup;
\end{vhdlcode}

\code{250 ps} is a value of type \code{time}, not an integer that you have
agreed to interpret as picoseconds. Adding it to a voltage is a compile error.
\SV{} has \code{time} as a 64-bit integer scaled by \code{timeunit} and
\code{timeprecision} directives that are per-file and easy to get wrong; the
number \code{250} in \code{\#250} means whatever the current time unit says it
means.

\subsection{User-defined resolution functions}
\index{resolution function}

\code{std\_logic} is not built into \VHDL{}. It is an ordinary enumerated type
in an ordinary package, and its multi-driver behaviour comes from an ordinary
function that you could have written. If you need a bus where multiple drivers
resolve by summing, or by taking the maximum, or by flagging a conflict as an
error, you write the resolution function and declare the signal to use it.

\SV{} has exactly one resolution rule, wired into the language for \code{wire}
and friends, and no way to change it. In practice this matters less than it
sounds, because almost nobody writes custom resolution functions. It is still
a real expressive difference, and it explains why \code{std\_logic} feels so
principled compared with \SV{}'s four-state \code{logic}.

\subsection{Operator overloading}
\index{operator overloading}

In \VHDL{} you can define \code{"+"} for your own type. This is why
\code{numeric\_std} exists at all, why \code{fixed\_pkg} can write \code{a * b}
for fixed-point values, and why a package of complex numbers is twenty lines
long. \SV{} has no operator overloading of any kind. A complex multiply is a
function call, and it looks like one.

\subsection{Configurations}
\index{configuration}

A \VHDL{} configuration is a separate design unit that binds entities to
architectures and component instances to entities, without touching the
design. It is how you swap a behavioural model for a gate-level netlist for
one instance in one simulation.

\SV{} has \code{config} blocks in the \LRM{}, and almost no support for them
in tools. In practice you do the same job with file lists,
\chOneDir{ifdef}, or a wrapper module, and you lose the compile-time record of
what was bound to what.

Note the asymmetry, because it is the fairest way to state this one. \SV{} has
a construct \VHDL{} lacks that points the other way: \code{bind} attaches an
extra module instance to an existing one from outside, without editing the
design, which is how assertions and coverage are kept out of synthesisable
files (\cref{ch:beyond}). So \VHDL{} can swap an implementation cleanly and
\SV{} can attach an observer cleanly, and neither can do the other's trick.

\subsection{Fixed- and floating-point packages}
\index{fixed\_pkg}\index{float\_pkg}

\VHDL{}-2008 ships \code{ieee.fixed\_pkg} and \code{ieee.float\_pkg} in the
standard library. These are not toys. \code{sfixed} and \code{ufixed} are
arrays indexed by a signed range where index 0 is the units bit, so the type
carries the binary point:

\begin{vhdlcode}
signal a : sfixed(2 downto -13);   -- Q3.13
signal b : sfixed(1 downto -14);   -- Q2.14
signal p : sfixed(4 downto -27);   -- exactly right for a*b
begin
  p <= a * b;                      -- width and binary point tracked
  q <= resize(p, q'high, q'low);   -- explicit, checked rounding
\end{vhdlcode}

The width of a product, the alignment of an addition, and the rounding of a
resize are all computed by the package from the types. The compiler catches
the mistake where you add a Q1.14 to a Q2.13 without shifting.

\SV{} has nothing comparable. There is no standard fixed-point library, no
binary-point tracking, and no compiler help. You write
\code{logic signed [15:0]}, you write \code{// Q1.14} in a comment, and you
keep the alignment right in your head:

\begin{svcode}
// a is Q1.14, b is Q1.14  ->  product is Q2.28 in 32 bits
logic signed [31:0] prod;
assign prod = a * b;
// take Q1.14 back out: >>> 14, then truncate. Rounding is your problem.
assign result = prod[29:14];
\end{svcode}

This matters directly for \refexamplebook. The CORDIC rotator is fixed-point
arithmetic from end to end, and every shift, every truncation and every
saturation is checked by nothing but review. A \VHDL{} implementation of the
same algorithm would have caught several classes of error at compile time. Be
disciplined: put the format in the type name (\code{data\_t} for Q1.14 in
\code{cordic\_pkg}), never mix formats without a comment saying why, and write
assertions on the numeric range.

\subsection{Deferred constants and separate package bodies}

A \VHDL{} package declares what is visible; the package body holds the
implementation, and can be recompiled without recompiling everything that uses
the package. A deferred constant lets the value itself live in the body, so a
change to a lookup table does not force a recompilation of every client. \SV{}
packages have no body: everything is in one file, and every change invalidates
every dependent compilation unit. On a large design this is a real difference
in build time.

\subsection{Protected types, generic packages, and 2019}

\VHDL{}-2002 protected types give encapsulated shared state with mutual
exclusion, which is most of what a \SV{} class does (\refverificationbook) minus
inheritance.
\VHDL{}-2008 generic packages and generic types give a genuine template
mechanism, closer in spirit to a C++ template than \SV{} parameters are. Both
are well specified. Support is the problem: generic packages work in the major
simulators, thinly in synthesis; \VHDL{}-2019 interfaces, mode views and
conditional analysis are excellent ideas with, at present, one or two
implementations between them.

\subsection{What you can do with a zero tool budget}
\index{open-source tools}

This one is decided by which half of the language a free tool has to cover.

For \VHDL{}, free simulation is mature and covers essentially the whole
language, \VHDL{}-93, -2002 and -2008 alike, at a speed that is perfectly
usable for a block-level regression. A student with no licences can analyse,
elaborate, simulate and dump waveforms for any \VHDL{} design in this book,
and the language they are learning is the whole language.

For \SV{} the free tools implement a subset, and the subset is not an
arbitrary one. \tool{Verilator} is excellent, \refsimulationbook is devoted
to it, and it is a compiler rather than an event-driven kernel. It does not
read \VHDL{} at all, so a mixed-language hierarchy needs a commercial
simulator. It is two-state, so the reset bug whose symptom is an \code{x}
propagating through the design does not appear. It implements neither
covergroups nor the constraint solver, and it has no route to UVM.
\refsimulationbook states those four boundaries precisely and gives the feature-by-feature position for the
version used in this book. \tool{Icarus Verilog} covers Verilog-2005 well and
\SV{} patchily, and it is not a tool to build a flow on.

Put the two together and the asymmetry is uncomfortable. If your budget is
zero and your language is \VHDL{}, you can do everything. If your language is
\SV{}, you can do the design half completely and quickly, and what falls
outside the free tools is exactly the constrained-random and coverage
machinery of \cref{sec:philosophy:verif}, which is the part this chapter has
just spent ten pages calling the reason to move.

\subsection{Verification frameworks that need no licence either}
\index{OSVVM}\index{UVVM}

\cref{sec:philosophy:verif} already said that OSVVM and UVVM give \VHDL{} most
of the same verification capability. Two points belong here rather than there,
because they are about the comparison rather than about the technique.

The first is that OSVVM's randomisation is coverage-driven: it biases the next
draw towards bins that are still empty, instead of drawing uniformly and
hoping. For many problems that closes the loop faster than a raw constraint
solver does, and it is a genuinely better idea, not a poor imitation of one.

The second is the learning cost. A working OSVVM testbench is a day's work. A
working UVM testbench is not, as \refverificationbook will make plain. Against that,
what \SV{} has and they do not is the constraint solver as a language feature,
assertions as a language feature, and the industry's momentum: the
verification IP you buy comes as UVM, and the job advertisements ask for UVM.

% ===========================================================================
\section{Which language for which job}
\label{sec:philosophy:which}
\label{sec:beyond:which}

The previous section listed what \VHDL{} does better; the rest of the chapter
listed what \SV{} does better. Neither list settles anything by itself,
because the question is never ``which language is better'' but ``which
language for this project''. So: no advocacy, criteria.

\textbf{Choose \VHDL{} for \RTL{} when} the design is arithmetic-heavy and
would benefit from \code{fixed\_pkg}; when the team is already fluent and the
cost of retraining is not repaid; when the flow is safety-critical and the
stronger type system is worth real money in review time; when the tool budget
is zero, since free \VHDL{} simulation covers the whole language and the free
\SV{} tools do not; or when the customer requires it, which in aerospace and
rail is common.

\textbf{Choose \SV{} for \RTL{} when} the design is control-heavy and
structural, with many similar blocks and wide buses that interfaces and packed
structs make tractable, both of which \cref{ch:beyond} covers; when you need parameterised
types across a hierarchy; when the same source must be read by a Verilog-only
tool, which is most open-source \RTL{} and most vendor IP; or when the
verification is going to be \SV{} anyway and you want one language in the
repository.

\textbf{Choose \SV{} for verification} in almost all cases where the
environment must scale: constrained random, functional coverage, assertions,
and UVM are a coherent package, and the industry has standardised on it. That
package is what \refverificationbook teaches. The exception is a
\VHDL{}-only shop with a working OSVVM or UVVM environment, and that exception
is legitimate.

\textbf{Mix them.} This is the common industrial answer and it works.
Practically every commercial simulator supports a mixed-language hierarchy:
\VHDL{} \RTL{} instantiated inside a \SV{} testbench, or the reverse. In
\tool{QuestaSim} you compile each language with its own compiler,
\code{vcom -2008} and \code{vlog -sv}, into the same library, and elaborate
the two together. The seams are the port types (map
\code{std\_logic\_vector} to packed vectors and keep the boundary plain: no
records, no interfaces, no enumerations across the boundary) and the inability
to bind \SV{} assertions into a \VHDL{} instance. Plan the boundary
deliberately and the arrangement is comfortable. Note that free tooling is not
an option here at all, because \tool{Verilator} does not read \VHDL{}.

\begin{notebox}{What actually decides it}
In practice the choice is made by the existing code base, the customer, and
the tool licences, in that order. Language merit comes fourth. Knowing both
languages well is worth more than a correct opinion about which is better.
\end{notebox}


% ===========================================================================
\section{Checklist}
\label{sec:philosophy:checklist}
% ===========================================================================

\begin{itemize}
\item \VHDL{} was designed as a language and grew a synthesis subset. \SV{}
  was designed as a simulator input and grew a type system. Every difference
  in this chapter follows from that.
\item Correctness in \VHDL{} comes from the analyser; in \SV{} it comes from
  the linter, the coding standard and simulation. Wire
  \code{verilator --lint-only -Wall} into your build on day one.
\item Widths never have to match in \SV{}. Assignments, expressions and port
  connections truncate or extend silently. Check every one.
\item \code{logic} is \code{std\_logic} without strengths and with a single
  driver. Use it for all RTL. Do not use \kw{int} or \kw{bit} for hardware:
  two-state types start at zero and hide missing resets.
\item Put \chOneDir{default\_nettype none} at the top of every file and
  \chOneDir{default\_nettype wire} at the bottom, or an undeclared name will
  silently become a one-bit wire.
\item When you see \code{library ieee; use ieee.numeric\_std.all;}, write
  \code{import my\_pkg::*;}. Put every shared \code{typedef} and
  \code{parameter} in a package; never declare anything at compilation-unit
  scope.
\item \chOneDir{define} and \chOneDir{include} are textual and unscoped. Prefix
  macro names, guard included files, and prefer a package parameter whenever
  one will do.
\item \VHDL{} signal assignment $\equiv$ \SV{} \code{<=}; \VHDL{} variable
  assignment $\equiv$ \SV{} \code{=}. \code{<=} evaluates in Active and updates
  in NBA: \VHDL{}'s evaluate/update pair, recreated by hand.
\item Non-blocking for everything a clock edge writes, blocking for everything
  combinational, never both in one block, never two blocks writing the same
  variable. \VHDL{} gave you this for free; \SV{} does not.
\item Declare intent: \code{always\_ff} for sequential, \code{always\_comb} for
  combinational, \code{always\_latch} when you mean it. One clock and at most
  one asynchronous reset per \code{always\_ff}.
\item In a testbench, drive and sample synchronous signals through a clocking
  block. \code{input \#1step} samples in Preponed, before any design code has
  run in that time slot.
\item Verification in \SV{} is a software problem: transactions, constrained
  random stimulus, functional coverage as the exit criterion, classes for
  reuse. Write the checker before the generator: random stimulus without a
  scoreboard proves nothing.
\item OSVVM and UVVM give \VHDL{} most of the same capability, and OSVVM's
  coverage-driven randomisation closes the loop faster than a raw constraint
  solver on many problems. The difference is ecosystem, industrial default and
  learning cost, not possibility.
\item Know what you are giving up, because it is not nothing: compile-time
  width and signedness checking, genuinely distinct user-defined types,
  physical types with units, user-defined resolution, operator overloading,
  \code{fixed\_pkg} and \code{float\_pkg}, configurations, deferred constants
  and separate package bodies. Compensate with lint, naming and review, in that
  order.
\item There is no \SV{} fixed-point library. Put the format in the type name,
  never mix formats without a comment saying why, and assert the numeric
  range. This is the single largest source of unchecked error in the CORDIC of
  \refexamplebook.
\item On a zero tool budget, \VHDL{} gets the whole language and \SV{} gets a
  fast subset: no \VHDL{} at the boundary, two-state only, no covergroups, no
  constraint solver, no UVM. Plan the flow around that, not around a hope.
\item The language choice is made by the existing code base, the customer and
  the tool licences, in that order, and by language merit fourth. Mixed-language
  hierarchies are normal and need a plain port boundary.
\end{itemize}
